\chapter{Selected Approach: Framework Design, Implementation, and Experimental Evaluation}
\chaptermark{Selected Approach}
\label{chap:selected-approach}

Chapter~\ref{chap:introduction} defined the need for a controlled and
traceable way of exercising healthcare-interoperability workflows.
Chapter~\ref{chap:stateofart} then established the methodological
requirements that follow from this need: the workload must preserve semantic
dependencies, distinguish request- and flow-level observations, expose its
temporal model, and keep client-side evidence separate from optional
middleware-side measurements. This chapter presents the selected approach from
the inside and records the evidence that is currently available.

The implementation described here is the Request Dataset Generator (RDG)
repository at the revision inspected for this migration. The chapter
distinguishes three claims that would otherwise be easy to conflate. A feature
is \emph{implemented} when it is present in the current code; it is
\emph{functionally exercised} when a test or archived run covers it; and it is
\emph{experimentally evaluated} only when an identified configuration,
execution log, normalised dataset, and metric report support the corresponding
analysis. In particular, mock execution validates the framework pipeline but
does not measure the capacity of a live regional or industrial service.

\section{Technical Objectives and System Architecture}
\label{sec:rdg-objectives-architecture}

\paragraph{Technical objectives.}
The framework translates the methodological requirements of
Section~\ref{sec:framework-implications} into five technical objectives. First,
it must generate synthetic healthcare data without requiring personal data.
Second, it must compose those data into ordered multi-step flows whose
dependencies remain explicit. Third, it must schedule independent flow
occurrences according to a configurable temporal model while preserving the
order of steps within each occurrence. Fourth, it must emit structured evidence
at request, flow, and time-window granularity. Finally, it must make
configuration, random seeds, execution mode, and output provenance recoverable
from each run.

These objectives delimit the artefact. RDG is not a conformance-certification
tool, a production audit system, or a clinical simulator. It creates and
executes controlled workloads at the client boundary and can enrich the
resulting records when a compatible middleware log is supplied. The absence of
that second source remains visible in the normalised dataset rather than being
replaced by estimated values.

\paragraph{End-to-end pipeline.}
The command dispatcher in \texttt{main.py} exposes independent commands for
patient generation, flow-dataset construction, timeline production, execution,
log collection, and metric calculation. The \texttt{pipeline} command connects
the same stages in memory and on disk:
\begin{align*}
  \text{configuration}
  &\rightarrow \text{synthetic data and flows}
  \rightarrow \text{timeline}\\
  &\rightarrow \text{execution log}
  \rightarrow \text{normalised dataset}
  \rightarrow \text{metrics and reports}.
\end{align*}
This separation allows an intermediate artefact to be inspected or replayed
without changing the responsibilities of the surrounding stages.

\begin{table}[htbp]
    \centering
    \small
    \begin{tabular}{|p{2.0cm}|p{3.7cm}|p{6.3cm}|}
        \hline
        \textbf{Stage} & \textbf{Primary implementation} & \textbf{Produced or transformed evidence} \\
        \hline
        Data generation & \texttt{fhir\_patient\_}\newline\texttt{generator.py} & Synthetic Patient resources and transaction Bundles. \\
        \hline
        Workload composition & \texttt{request\_dataset\_}\allowbreak\texttt{generator.py}, \texttt{flow\_}\allowbreak\texttt{orchestrator.py} & Weighted population of flows, ordered steps, identifiers, dependencies, and payload metadata. \\
        \hline
        Temporal scheduling & \texttt{timeline\_}\newline\texttt{generator.py} & Timestamped flow occurrences with base/burst provenance and experiment identifiers. \\
        \hline
        Execution & \texttt{traffic\_exec\_}\newline\texttt{engine.py} & Mock or live step outcomes and JSONL request/response/flow events. \\
        \hline
        Collection & \texttt{log\_collector.py} & Request, flow, and bucketed time-series populations; optional middleware enrichment. \\
        \hline
        Analysis & \texttt{metric\_engine.py} & Machine-readable metrics, human-readable report, CSV tables, warnings, and optional graphs. \\
        \hline
    \end{tabular}
    \caption{Verified responsibilities of the RDG pipeline.}
    \label{tab:rdg-pipeline}
\end{table}

\paragraph{Repository organisation and dependency direction.}
The top-level modules implement the pipeline stages, while the transaction
package contains the transaction builders and the flow
orchestrator. A transaction builder returns one executable step with its
method, target selected from configuration, headers, body, transport profile,
dependency specification, output-extraction rules, and---in mock mode---a
synthetic response. The orchestrator adds correlation metadata and combines
these steps into flow types. The execution engine consumes this common shape;
it does not contain healthcare-specific request templates.

This direction of dependency is important. Domain-specific contracts remain in
the transaction modules, composition remains in the orchestrator, and temporal
or analytical mechanisms operate on shared metadata. Adding a new transaction
therefore does not require the scheduler or metric engine to understand its
payload, provided that the module satisfies the common step contract and
declares the outputs required by downstream steps.

\paragraph{Configuration and run provenance.}
Configuration is loaded from JSON and merged with command defaults. Before a
pipeline execution, the dispatcher establishes a scenario identifier, run
identifier, execution identifier, and seeds for workload, timeline, and fault
injection. It also resolves output paths under a run-specific directory and
writes the effective configuration as \texttt{config.resolved.json}. The run
directory can then contain the generated dataset, timeline, JSONL execution
log, normalised metrics, metric reports, tabular summaries, and graphs.

The effective configuration is more authoritative than the current scenario
template when interpreting an archived run: templates may evolve after an
execution, whereas the resolved copy records the parameters attached to that
run. Credentials, security tokens, certificate material, private endpoints,
and generated clinical-like values are deliberately excluded from this
chapter. Reproducibility concerns non-sensitive parameters and provenance, not
publication of protected configuration.

\section{Synthetic Inputs and Executable Transactions}
\label{sec:rdg-data-workload-transactions}

\paragraph{Synthetic Patient resources.}
The data generator creates a FHIR Patient-shaped resource with synthetic
identifiers, demographics, contact information, addresses, regional
extensions, a contained practitioner, and a contained assistance contract.
It can also construct the transaction Bundle used by the RVE-55 path and
serialise generated structures as JSON or XML. Random generation removes the
need for real patient records, while an optional identifier and a declared seed
support controlled regeneration.

The generated structure is intended to exercise the contracts implemented by
the framework. Its richness does not prove conformance with every cardinality,
terminology rule, or business constraint of the applicable regional profile.
That distinction mirrors the boundary established in
Section~\ref{subsec:fhir-r4}: the code demonstrates what the workload can
produce, while the primary specification remains authoritative for
conformance.

\paragraph{Workload population and correlation model.}
\texttt{build\_workload} accepts the number of source flows, a weighted
\texttt{flow\_mix}, an execution mode, the complete configuration, and an
optional seed. Weights are normalised and sampled to choose a flow type; a new
synthetic patient is generated for each source flow. The resulting flow carries
a flow identifier, a type, and an ordered list of steps. Every step has its own
request identifier and index, while the timeline later adds a distinct
\texttt{flow\_occurrence\_id} whenever a source flow is scheduled. This last
identifier prevents repeated uses of the same source flow from being merged
during log reconstruction.

The current implementation defines the seven flow types in
Table~\ref{tab:implemented-flow-types}. The table reports software coverage,
not equal validation status. In particular, a zero default weight or presence
in the map does not establish that a path has been exercised against a live
provider.

\begin{table}[htbp]
    \centering
    \scriptsize
    \begin{tabular}{|p{3.2cm}|p{7.8cm}|p{1.0cm}|}
        \hline
        \textbf{Flow type} & \textbf{Ordered transaction modules} & \textbf{Steps} \\
        \hline
        \texttt{FLOW\_CONTEXT\_CALL} & RVE-1.b $\rightarrow$ RVE-121 $\rightarrow$ RVE-130 & 3 \\
        \hline
        \texttt{FLOW\_PATIENT\_QUERY} & RVE-1.b $\rightarrow$ RVE-54 & 2 \\
        \hline
        \texttt{FLOW\_PATIENT\_CREATE} & RVE-1.b $\rightarrow$ RVE-55 & 2 \\
        \hline
        \texttt{FLOW\_ITI\_18} & RVE-1.b $\rightarrow$ ITI-18 & 2 \\
        \hline
        \texttt{FLOW\_ITI\_43} & RVE-1.b $\rightarrow$ ITI-43 & 2 \\
        \hline
        \texttt{FLOW\_PATIENT\_QUERY\_}\allowbreak\texttt{MIDDLEWARE} & RVE token $\rightarrow$ gateway RVE-54 $\rightarrow$ gateway RVE-55 $\rightarrow$ gateway RVE-57 $\rightarrow$ gateway RVE-100 & 5 \\
        \hline
        \texttt{FLOW\_SCRYBASIGN\_}\allowbreak\texttt{SIGN} & GetUserInfo $\rightarrow$ SignOneDoc & 2 \\
        \hline
    \end{tabular}
    \caption{Flow types present in the current orchestrator map.}
    \label{tab:implemented-flow-types}
\end{table}

\paragraph{Step dependencies and the mock/live boundary.}
A flow is more than a list of independent HTTP requests. A step may declare
that a field comes from an output of an earlier step. In live mode, downstream
headers, bodies, and URLs can contain scoped placeholders; the execution engine
resolves them from the per-occurrence context only after the preceding response
has been processed. Required outputs are validated before the chain continues.
An unresolved placeholder or missing required output is logged as a failure
instead of sending an incomplete dependent request.

In mock mode, transaction builders attach synthetic responses and the
orchestrator propagates their declared outputs while constructing the chain.
Execution then consumes those responses without opening an HTTP session. This
mode tests composition, scheduling, failure handling, logging, collection, and
metric production under controlled conditions. It cannot measure network,
gateway, provider, or production-system behaviour. Dry-run is a separate
execution branch that records a successful empty interaction without sending
requests and must not be interpreted as either mock-contract validation or live
service evidence.

\paragraph{RVE, ITI, and mediated paths.}
The direct regional paths combine an initial assertion-producing transaction
with context, patient-registry, or document transactions. The two ITI flows use
dedicated transaction modules after the initial integration step. The mediated
patient path has a different structure: it obtains a middleware token and then
executes the configured gateway scenarios for RVE-54, RVE-55, RVE-57, and
RVE-100. This five-step path is therefore not a request-for-request replacement
for the two-step direct patient query. Comparisons between their completion
times must acknowledge the different step count and contracts.

The code confirms that RVE-57 is present in the mediated flow. It does not, by
itself, provide the primary domain evidence required to finalise the semantic
description left open in Section~\ref{sec:interoperability-transactions}.
[SOURCE NEEDED: authoritative RVE-57 specification and approved mapping of the
gateway scenarios to the thesis use cases.]

\paragraph{Digital-signature provider flow.}
\texttt{FLOW\_SCRYBASIGN\_SIGN} composes two SOAP POST steps in a fixed order:
the provider user-information operation and the single-document signing
operation. Both builders obtain authentication and mutual-TLS profile
information through scoped configuration helpers. Their mock branches return a
minimal successful SOAP envelope, which allows the generic pipeline to
exercise the two-step shape without contacting the provider.

The current signing request builder contains a fixed test document and does not
declare a data dependency between the first and second response. This is the
implemented contract and should not be described as a dynamically generated
clinical-document signing workflow. Unit tests verify the order, SOAP operation
shape, authentication-header construction, and transport-profile selection,
but no archived run in the available corpus contains this flow type.
The newer sanitised live corpus supersedes the preceding archive statement. It
records 17{,}838 HTTP~200 successes for
\texttt{SCRYBASIGN-GET-USER-INFO} and 17{,}838
\texttt{SCRYBASIGN-SIGN-ONE-DOC} outcomes, of which 17{,}015 are HTTP~200
successes and 823 are failures. Of the failures, 28 have no HTTP response and
795 have status HTTP~502 with declared fault injection. These observations
establish live execution and client-side outcome coverage, but not provider-side
response semantics or authorisation of the test.

[DATA MISSING: approved mock validation, confirmation that the live execution
was credential-authorised, sanitised provider-side evidence, and validation of
the signed-response semantics.]

\section{Traffic Execution and Evidence Pipeline}
\label{sec:rdg-execution-observability}

\paragraph{Temporal model and occurrence generation.}
The timeline generator separates flow content from arrival time. It supports a
constant rate, a configurable step profile, and a synthetic time-banded
clinical profile. For a variable intensity $\lambda(t)$, arrivals are produced
through Lewis--Shedler thinning: candidates are sampled from an upper-bound
homogeneous process and retained according to
$\lambda(t)/\lambda_{\max}$. Optional bursts add a configurable number of
closely spaced occurrences around selected baseline arrivals. Every event
records whether it is a base or burst arrival and, where applicable, its burst
identifier.

The generator samples from the already constructed source dataset according to
its empirical flow-type distribution. It therefore changes the temporal
population without modifying transaction definitions. A seed makes workload
and timeline sampling repeatable for a fixed code and configuration, but it
cannot make the latency of live remote services deterministic. The nominal
duration and rate also do not fix the final number of arrivals, which remains a
realisation of the stochastic process.

\paragraph{Scheduling and concurrency semantics.}
The execution engine follows scheduled timestamps using an asynchronous task
per flow occurrence. A semaphore limits concurrent flows and encloses the whole
flow rather than a single request. Inside that boundary, steps are executed in
order and share only their occurrence-local output context; separate flows can
overlap. The engine stops a chain at the first non-success status or processing
error, records the failing step, and leaves the remaining downstream steps
unexecuted.

Live execution opens an HTTP client session, resolves dependencies, optionally
uses the configured proxy and mutual-TLS context, sends the request, extracts
declared outputs, and validates required fields. Mock execution reads the
synthetic response attached to each step. A deterministic fault injector can
apply declared error, timeout, or latency-spike rules by transaction, flow
type, or step. These injected outcomes evaluate framework behaviour under
controlled faults; they are not observations of provider reliability.

\paragraph{Structured execution events.}
The execution logger emits three JSONL event types. A
\texttt{request\_sent} event records experiment identifiers, flow occurrence,
request and step identifiers, transaction, payload size, scheduled and actual
start times, queue delay, active flows and steps, and burst provenance. A
\texttt{response\_received} event records status, response size, client
latency, success, error, and fault-injection provenance. A
\texttt{flow\_completed} event records total duration, success, first failed
step, declared step count, and burst provenance.

This schema implements the distinction between an interaction and an
application flow established in Chapter~\ref{chap:stateofart}. It also
records enough provenance to stratify observations by scenario, flow,
transaction, arrival type, or execution. Payload contents, assertions, tokens,
credentials, and certificate data are not required by the analytical schema
and must not be exported as thesis evidence.

\paragraph{Normalisation and optional grey-box enrichment.}
The collector parses the framework JSONL stream, pairs sent and received events
by occurrence, reconstructs completed flows, and aggregates a bucketed time
series. Its output contains four top-level parts: metadata, requests, flows, and
timeseries. Request records retain client latency, outcome, concurrency,
scheduling, size, and provenance. Flow records retain completion time, executed
steps, skipped downstream steps, and outcome. Metadata records the source files
and the observed time range.

A second parser can ingest a compatible middleware JSONL stream. Matching
records may add gateway and upstream latencies, retry count, circuit-breaker
state, and infrastructure observations. When no middleware record is
available, these fields remain null. The present archived runs report zero
middleware events, so they support only the framework level of the grey-box
model. No gateway/upstream decomposition or infrastructure correlation can be
derived from them.

\paragraph{Metric and reporting engine.}
The metric engine calculates eighteen measures grouped into performance,
reliability, resilience, scalability, and observability. The performance group
includes flow completion time, step latency, request and flow throughput,
burst/\allowbreak non-burst summaries, and an optional direct-versus-mediated comparison.
Reliability includes error, timeout, retry, and flow-success rates. Resilience
includes recovery episodes, failure propagation, and availability over
time-series buckets. Scalability includes available resource observations,
performance versus load, client-observed concurrency, and throughput
efficiency. Observability includes tracing completeness, event correlation, and
scheduling fidelity.

Descriptive statistics include the population size, mean, median, standard
deviation, minimum, maximum, and selected percentiles. The engine can produce a
JSON report, a Markdown report, CSV summary tables, warnings, and graphs. Its
warning detector compares metrics with configurable values, including built-in
fallbacks. These are software warning thresholds, not domain-approved
acceptance criteria. They cannot close the gap identified in
Section~\ref{sec:approach-selection}.

\section{Evaluation Protocol and Experimental Design}
\label{sec:rdg-evaluation-methodology}

\paragraph{Evaluation questions.}
The implemented architecture and the evidence required by the problem
statement lead to four evaluation questions:
\begin{enumerate}
    \item Can RDG preserve ordered transaction dependencies and isolate the
    state of concurrent flow occurrences in both mock and live modes?
    \item Can a declared flow mix and temporal model be reproduced with
    traceable identifiers, seeds, configuration, and artefacts?
    \item Can execution events be transformed into consistent request-, flow-,
    and time-window populations without filling unavailable middleware fields?
    \item Within an approved scenario, how do latency distributions,
    throughput, reliability, queueing, and flow completion change with flow
    composition, offered load, bursts, or injected faults?
\end{enumerate}
The first three questions concern functional and methodological validity. The
fourth concerns experimental behaviour and must be answered separately for
each execution mode and observed system boundary.

\paragraph{Critical evaluation of the architecture.}
The pipeline provides strong separation between domain contracts, temporal
load, execution, and analysis. It also makes intermediate artefacts inspectable
and allows the same collector and metric engine to process mock and live logs.
The main architectural trade-off is that a generic step contract can preserve
ordering and metadata without establishing semantic equivalence among flows.
A two-step direct patient query and a five-step mediated path can both be
executed, but their difference cannot be attributed solely to middleware.

The execution semaphore applies to complete flows. This protects
within-occurrence dependencies and makes the number of active flows explicit,
but it also means that long or blocked flows retain a concurrency slot. This is
part of the evaluated execution model and should not be mistaken for an
unqualified model of provider-side concurrency. Similarly, the client-side
clock measures scheduling and request completion at the generator; internal
service queues remain outside that boundary unless a verified second source is
available.

\paragraph{Two-level grey-box measurement model.}
Level~1 is always generated by RDG and observes scheduled dispatch, actual
dispatch, client-side step latency, response outcome, flow completion, and
client-observed concurrency. Level~2 is optional and may observe middleware
internal intervals, retries, circuit-breaker state, or infrastructure
measurements. A join is acceptable only when identifiers and occurrence
semantics are compatible. The completeness of that join is a data-quality
measure, not an assumption.

The two levels must not be combined by visually aligning unsynchronised series
or subtracting intervals with different clock and boundary semantics. A
statistical association may support a diagnostic hypothesis, but causal
attribution requires controlled variation and compatible measurements. The
available run corpus contains Level~1 only.

\paragraph{Variables and controlled factors.}
The independent variables exposed by the framework include execution mode,
flow mix, source-dataset size, temporal profile and intensity, burst parameters,
concurrency limit, timeout, and declared fault-injection rules. Dependent
variables include step and flow latency distributions, request and flow
throughput, error and timeout rates, flow-success rate, scheduling drift,
queue delay, active flows and steps, recovery episodes, and evidence
completeness.

At minimum, a comparison must hold constant the code revision, transaction
modules, execution mode, credentials and transport policy where applicable,
flow mix unless it is the factor under study, seed policy, duration, concurrency
limit, metric implementation, and observation boundary. Payload-size and
step-count distributions must be reported when they differ across compared
flows. Live tests additionally require a documented hardware, operating-system,
network, remote-environment, and time-of-execution context.

\paragraph{Reproducibility protocol and missing controls.}
An analysable run must retain the code revision, resolved configuration,
dataset, timeline, execution log, normalised dataset, metric report, and any
excluded or failed-run record. Scenario, run, and execution identifiers must
agree across these artefacts. Seeds must be recorded separately for workload,
timeline, and fault injection. Results must identify the execution mode and
must never combine mock and live samples in one distribution.

The current archive satisfies much of the artefact chain for individual runs,
but it does not constitute the final approved experiment campaign. The
sanitised batch planned 30 repetitions for each of ten scenarios, whereas the
available complete population ranges from five to fifteen runs per scenario
and totals 74 runs. Provenance identifies source commit
\texttt{f08fb53fd4ad11ef803642c0a066e9a6d9e2b327} as the RDG checkout observed
at extraction time, rather than as a revision embedded in the batch manifest.

[DATA MISSING: approved warm-up and cool-down policy, minimum populations,
hardware and operating-system inventory, network path and
clock-synchronisation evidence, remote-environment version, run exclusion
policy, and a manifest-embedded immutable code/configuration revision.]

\section{Executed Scenarios, KPIs, and Quality Controls}
\label{sec:rdg-scenarios-kpis-quality}

\paragraph{Declared scenarios versus executed evidence.}
The repository contains configuration templates for constant baselines,
step-load ramps, bursts, direct-versus-mediated comparisons, fault injection,
flow-type comparisons, a small live smoke test, and an extreme stress profile.
Their presence describes intended experiment capabilities. Only a scenario
with a run directory and a coherent artefact chain is treated here as executed.
The resolved configuration in each run directory determines its mode and
parameters.

Table~\ref{tab:archived-rdg-runs} inventories the six available run
directories. Five are explicitly mock executions. The sixth is recorded as a
live, three-flow smoke run, but its size and coverage make it suitable only as
limited functional evidence.

\begin{table}[htbp]
    \centering
    \scriptsize
    \begin{tabular}{|p{3.0cm}|p{1.1cm}|p{1.0cm}|p{0.9cm}|p{5.1cm}|}
        \hline
        \textbf{Archived scenario} & \textbf{Mode} & \textbf{Flows} & \textbf{Req.} & \textbf{Observed flow types} \\
        \hline
        \texttt{baseline\_mock} & mock & 12 & 26 & context, patient query, patient create \\
        \hline
        \texttt{burst\_mock} & mock & 136 & 321 & context, patient query, patient create \\
        \hline
        \texttt{direct\_vs\_}\allowbreak\texttt{middleware\_mock} & mock & 76 & 287 & direct patient query, mediated patient path \\
        \hline
        \texttt{error\_injection\_mock} & mock & 65 & 151 & context, patient query, patient create \\
        \hline
        \texttt{flow\_type\_}\allowbreak\texttt{comparison\_}\allowbreak\texttt{middleware} & mock & 45 & 225 & mediated patient path only \\
        \hline
        \texttt{live\_smoke\_small} & live & 3 & 6 & direct patient query only \\
        \hline
    \end{tabular}
    \caption{Executed runs supported by resolved configurations and normalised datasets.}
    \label{tab:archived-rdg-runs}
\end{table}

The executed corpus does not cover ITI-18, ITI-43, or
\texttt{FLOW\_SCRYBASIGN\_}\allowbreak\texttt{SIGN}; nor does it include a step-ramp or extreme
stress run. Current scenario templates that mention these flows cannot be used
as evidence that the corresponding test was performed. The newer sanitised
live corpus supersedes this six-run inventory: it contains five complete
\texttt{load\_ramp} runs, five \texttt{stress\_extreme} runs, six
\texttt{flow\_type\_comparison\_direct} runs, fifteen
\texttt{flow\_type\_comparison\_middleware} runs, and live outcomes for
ITI-18, ITI-43, and both ScrybaSign transactions.

[EXPERIMENTAL RESULT MISSING: 226 of the 300 planned runs were not complete;
finish the planned repetitions before treating the campaign as final.]

[TO EXPAND: reconcile Table~\ref{tab:archived-rdg-runs} and the superseded
six-run inventory with the newer sanitised campaign in a later drafting phase.]

\paragraph{Operational KPI populations.}
Flow completion time is calculated over completed flow records and must be
stratified by flow type, outcome, and mode. Step latency is calculated over
paired request/response occurrences and is stratified at least by transaction
and flow type. Request throughput counts completed requests per observed second;
flow throughput counts completed flow occurrences over the same run interval.
The two are reported together because flow types have different step counts.

Reliability uses request errors, timeouts, and completed-flow outcomes. Failure
at a step is associated with skipped downstream steps rather than fabricated
requests. Scheduling fidelity uses the difference between planned and actual
client dispatch together with queue delay and active-flow/step counts.
Burst analysis uses the explicit base/burst labels. Middleware breakdown,
resource saturation, and retry or circuit-breaker measures are valid only for
records actually enriched by a second source.

\paragraph{Data-quality controls.}
Before interpretation, each run is checked for:
\begin{itemize}
    \item an available resolved configuration and explicit execution mode;
    \item consistent scenario, run, execution, flow-occurrence, flow, request,
    and step identifiers;
    \item paired sent/received events and an explicit flow-completion event;
    \item recorded population sizes, time range, units, and flow-type coverage;
    \item explicit fault-injection and burst provenance;
    \item separate mock and live populations;
    \item the count and coverage of middleware-enriched records;
    \item retained failed or incomplete observations rather than silent
    imputation.
\end{itemize}

All six normalised datasets identify their framework log, and their request
records are associated with reconstructed flows. They also all report zero
middleware events and zero requests with middleware timing fields. This is a
valid Level~1 dataset condition, but it prevents a Level~2 grey-box analysis.
[DATA MISSING: verified middleware log with compatible occurrence keys,
measurement-boundary documentation, clock evidence, and quantified join
coverage.]

\paragraph{Acceptance criteria.}
The archived reports contain configurable warning values and the engine
produces warnings when those values are crossed. They are useful for automated
inspection but are not clinical, contractual, or production-service acceptance
limits. The experiment can currently judge internal consistency, functional
completion, reproducibility of mock inputs, and descriptive behaviour within a
scenario. It cannot judge production suitability against an approved service
level.

[DATA MISSING: domain-supported acceptance thresholds, required confidence or
uncertainty, minimum sample sizes, repetition count, comparison rules, and the
scenario--KPI--population--threshold--evidence matrix defined ex ante.]

\section{Evidence, Discussion, and Design Guidelines}
\label{sec:rdg-results-discussion}

\paragraph{Functional evidence.}
The code-level test suite exercises placeholder resolution, supported flow
composition, ITI profiles, ScrybaSign request construction and ordering,
timeline occurrence identity, burst tagging, configuration validity, and
pipeline command dispatch. During this migration, 29 tests were executed in a
read-only environment: 28 completed successfully and the graph/table-output
test stopped because the optional plotting dependency was not installed. This
environmental failure does not invalidate the underlying metric calculation,
but it means that the current test execution is not a clean all-green release
record.

The archived mock runs further validate that generated flows can traverse the
timeline, executor, logger, collector, and metric engine. The small live run
records three successful direct patient-query flows, six successful requests,
and the expected two-step transaction shape. Its population is too small and
its environmental context too incomplete for performance inference; it is
reported only as limited live-path functional evidence.

\paragraph{Mock pipeline results.}
Table~\ref{tab:mock-rdg-results} reports selected outcomes from the five mock
runs. Values describe execution of synthetic responses and framework-side
processing on the recorded host. They must not be read as regional-service,
gateway, middleware, or provider latency.

\begin{table}[htbp]
    \centering
    \scriptsize
    \begin{tabular}{|p{3.4cm}|p{1.45cm}|p{1.45cm}|p{1.45cm}|p{1.45cm}|p{1.45cm}|}
        \hline
        \textbf{Scenario} & \textbf{FCT P50 ms} & \textbf{FCT P99 ms} & \textbf{Req/s} & \textbf{Error \%} & \textbf{Flow success \%} \\
        \hline
        Baseline mock & 0.462 & 0.859 & 11.770 & 0.000 & 100.000 \\
        \hline
        Burst mock & 0.694 & 2.134 & 27.821 & 0.000 & 100.000 \\
        \hline
        Direct/mediated mock & 1.796 & 3.033 & 18.162 & 0.000 & 100.000 \\
        \hline
        Fault-injection mock & 0.839 & 50.930 & 12.768 & 5.960 & 86.154 \\
        \hline
        Mediated-flow mock & 2.259 & 6.599 & 24.534 & 0.000 & 100.000 \\
        \hline
    \end{tabular}
    \caption{Selected framework-side outcomes from archived mock runs. FCT denotes flow completion time.}
    \label{tab:mock-rdg-results}
\end{table}

The baseline and burst runs show that burst-labelled occurrences are processed
and remain available for stratified analysis. The increase in reported P99 in
the burst run is an observation of this mock execution, not evidence of remote
system degradation. Only one seed and one execution are available, and the two
runs do not establish a distribution of run-to-run variability.

The direct-versus-mediated mock report contains 31 direct and 45 mediated flow
occurrences. Their mean completion times are respectively 0.781~ms and
2.166~ms, with P95 values of 1.147~ms and 2.961~ms. The report calls the
difference an \emph{apparent overhead}; this wording is essential. The paths
have two and five steps respectively, use synthetic responses, and have no
middleware-internal observations. The difference validates stratified
reporting but cannot isolate middleware cost.

The fault-injection run contains 65 flows and 151 executed requests. Nine flows
fail after declared synthetic faults, giving a flow-success rate of
86.154\% and a request error rate of 5.960\%. The report identifies two
recovery episodes with a mean duration of 3.5~s. These values demonstrate that
the pipeline retains injected failures, interrupts downstream steps, and
computes resilience summaries. They do not estimate the failure or recovery
behaviour of a live service.

\paragraph{Grey-box result boundary.}
Although the direct-versus-mediated and mediated-only runs exercise the
five-step gateway path, their normalised datasets contain no middleware events.
Consequently, gateway latency, upstream latency, retry count, circuit-breaker
state, CPU, memory, and internal connection measurements are unavailable. The
only defensible comparison is between client-observed flow shapes in mock mode.
The Level~2 portion of the planned grey-box evaluation remains unexecuted.

[EXPERIMENTAL RESULT MISSING: execute matched direct and mediated live
scenarios with a real, documented middleware log; verify occurrence-level
joins and observation boundaries before analysing internal latency or resource
relationships.]

\paragraph{Relation to the state of the art.}
The available evidence supports the methodological conclusions of
Chapter~\ref{chap:stateofart} at the level of the artefact. Step and flow
populations are both retained, percentiles expose distribution tails that a
mean would hide, and structured occurrence identifiers prevent repeated
source flows from being collapsed. The evidence does not yet support the
stronger empirical comparison with the reviewed literature: different systems,
payloads, network paths, and experimental designs preclude transfer of their
quantitative findings.

In particular, the mock tail and fault-injection observations illustrate why
P95/P99 and failed-flow populations are useful, but do not validate a
production scheduling or capacity rule. Likewise, the small live smoke run
does not establish that lower-level optimisations improve end-to-end
performance. Those questions require the missing repeated live campaign and
declared acceptance framework.

\paragraph{Design and evaluation guidelines.}
The implementation and current evidence support the following conservative
guidelines:
\begin{enumerate}
    \item preserve transaction, step, flow, and occurrence identities from
    generation through reporting;
    \item treat step order and dependency propagation as correctness
    invariants, while applying concurrency between independent flows;
    \item record resolved configuration and independent seeds with every run;
    \item separate mock, dry-run, and live evidence in storage and analysis;
    \item report request and flow throughput together because flow shapes have
    different step counts;
    \item qualify direct-versus-mediated comparisons by semantic equivalence,
    step count, payload, and observation boundary;
    \item leave unavailable middleware and infrastructure fields null and
    quantify their coverage;
    \item distinguish configurable software warnings from approved acceptance
    thresholds;
    \item calibrate flow mixes and arrival profiles before claiming
    representativeness of operational traffic;
    \item derive production or capacity guidance only from repeated,
    controlled, appropriately authorised live experiments.
\end{enumerate}

\paragraph{Residual limitations and contribution boundary.}
The present chapter can substantiate the architecture, current transaction map,
mock/live boundary, evidence pipeline, archived mock validation, and one
limited live patient-query smoke run. It cannot substantiate production
capacity, provider reliability, causal middleware attribution, operational
traffic representativeness, or the performance of ITI and ScrybaSign paths.
These limitations preserve the distinction between software capability and
experimental evidence established at the beginning of the chapter.

The newer sanitised live corpus provides repeated execution evidence for
ITI-18 (21{,}722 outcomes, including 3{,}395 successes), ITI-43
(21{,}514 outcomes, including 21{,}371 successes),
\texttt{SCRYBASIGN-GET-USER-INFO} (17{,}838 successes), and
\texttt{SCRYBASIGN-SIGN-ONE-DOC} (17{,}838 outcomes, including 17{,}015
successes). The planned protocol remains incomplete, and the sanitised corpus
contains no Level~2 middleware evidence.

[EXPERIMENTAL RESULT MISSING: complete the remaining planned repetitions and
integrate sanitised Level~2 middleware evidence before drawing final
cross-layer or capacity conclusions.]

[TO EXPAND: after the final campaign, discuss every accepted and rejected
hypothesis against the ex ante criteria, report uncertainty and failed runs,
and derive only the architectural, performance-profiling, and gateway-evaluation
guidelines supported by those results.]
